\clearpage
\thispagestyle{empty}\SourceLeaf{201}
\pdfbookmark[0]{General Theorems on Forces of Attraction and Repulsion That Vary Inversely with the Square of the Distance}{allgemeine-lehrsaetze-title}
\begin{center}
\vspace*{0.6in}
{\Large GENERAL THEOREMS\par}
\bigskip
ON FORCES OF ATTRACTION AND REPULSION\par
\bigskip
THAT VARY INVERSELY WITH\par
\bigskip
{\Large THE SQUARE OF THE DISTANCE\par}
\vspace{0.5in}
BY\par
\bigskip
{\large CARL FRIEDRICH GAUSS.\par}
\vfill
\rule{0.85\textwidth}{0.3pt}\par
Resultate aus den Beobachtungen des magnetischen Vereins im Jahre 1839.\par
Edited by \textsc{Gauss} and \textsc{Weber}. Leipzig 1840.\par
\rule{0.85\textwidth}{0.3pt}
\end{center}
\clearpage
\SourceLeaf{203}
\fancyhead[L]{\small General Theorems on Forces of Attraction and Repulsion}
\renewcommand{\Article}[1]{\par\medskip\Needspace{4\baselineskip}\pdfbookmark[1]{Art. #1}{lehrsaetze-art-#1}\centerline{\textbf{#1.}}\nobreak\smallskip}
\pdfbookmark[0]{General Theorems on Forces of Attraction and Repulsion}{lehrsaetze-body}
\begin{center}
GENERAL THEOREMS\\
ON FORCES OF ATTRACTION AND REPULSION\\
THAT VARY INVERSELY WITH\\
THE SQUARE OF THE DISTANCE.
\end{center}
\Article{1}
Nature shows us a variety of phenomena that we explain by assuming that the smallest particles of substances exert forces on one another, in inverse proportion to the squares of their mutual distances.

Foremost among these is universal gravitation Under its action, each ponderable molecule $\mu$ exerts a motive force on another $\mu'$. If the distance is set $=r$, that force is given by $\frac{\mu\mu'}{rr}$ and tends to bring them closer along the straight line joining them.

If two magnetic fluids are assumed in order to explain magnetic phenomena, with one regarded as a positive quantity and the other as negative, two such elements $\mu$, $\mu'$ also exert a motive force on each other, measured by $\frac{\mu\mu'}{rr}$. It acts along the straight line joining them, as a repulsion if $\mu$, $\mu'$ are of the same kind and as an attraction if they are of different kinds.

Much the same applies to the way the particles of the electric fluids act on one another.

\SourceLeaf{204}
The linear element $\mathrm{d}s$ of a galvanic current likewise exerts a motive force on an element $\mu$ of the magnetic fluid (if that fluid is admitted), inversely proportional to the square of the distance $r$. but here an altogether different feature also comes into play. The force is directed perpendicular to the plane laid through $\mu$ and the direction of $\mathrm{d}s$, rather than along the straight line joining the elements, and its magnitude also depends on the angle that $r$ makes with the direction of $\mathrm{d}s$, in addition to the distance. If this angle is called $\theta$, then $\frac{\sin\theta.\mu\mathrm{d}s}{rr}$ measures the motive force that $\mathrm{d}s$ exerts on $\mu$, and $\mu$ exerts an equally large force on the current element $\mathrm{d}s$ or on its ponderable carrier, in a direction parallel and opposite to that of the first force.

If, with \textsc{Ampère}, it is assumed that two elements $\mathrm{d}s$, $\mathrm{d}s'$ of galvanic currents attract or repel each other along the straight line joining them, the phenomena compel us to take this force too as inversely proportional to the square of the distance, and at the same time they require a somewhat more complicated dependence on the direction of the current elements.

In this treatise we will confine ourselves to the first three cases, or to forces that act along the straight line between the element exerting the force and the element it acts on and are simply inversely proportional to the square of the distance, although several theorems also apply to the other cases with slight changes. A detailed treatment of those cases must be reserved for another treatise.

\Article{2}
We use $a$, $b$, $c$ for the rectangular coordinates of a material point from which a repulsive or attractive force acts. The accelerating force itself at an arbitrary point $O$, whose coordinates are $x$, $y$, $z$, is denoted by
\[
\frac{\mu}{(a-x)^2+(b-y)^2+(c-z)^2}=\frac{\mu}{rr}
\]
Here $\mu$ denotes, in the first case of the preceding article, the ponderable matter
\SourceLeaf{205}
present at the first point, and in the second and third cases the quantity of magnetic or electric fluid. Resolving this force parallel to the three coordinate axes gives the components
\[
\frac{\varepsilon\mu(a-x)}{r^3},\qquad \frac{\varepsilon\mu(b-y)}{r^3},\qquad \frac{\varepsilon\mu(c-z)}{r^3}
\]
where $\varepsilon=+1$ or $=-1$ according as the force acts attractively or repulsively, as follows from the nature of what exerts the force and what receives it. These components take the form of the partial derivatives
\[
\frac{\mathrm{d}\frac{\varepsilon\mu}{r}}{\mathrm{d}x},\qquad \frac{\mathrm{d}\frac{\varepsilon\mu}{r}}{\mathrm{d}y},\qquad \frac{\mathrm{d}\frac{\varepsilon\mu}{r}}{\mathrm{d}z}
\]
Thus if the same point $O$ is acted on by several sources of force $\mu^0$, $\mu'$, $\mu''$ and so on from distances $r^0$, $r'$, $r''$ and so on, and if we set
\[
\frac{\mu^0}{r^0}+\frac{\mu'}{r'}+\frac{\mu''}{r''}+\text{ etc.}=\Sigma\frac{\mu}{r}=V
\]
the components of the total force acting at $O$ are given by
\[
\frac{\varepsilon\mathrm{d}V}{\mathrm{d}x},\qquad \frac{\varepsilon\mathrm{d}V}{\mathrm{d}y},\qquad \frac{\varepsilon\mathrm{d}V}{\mathrm{d}z}
\]

When the sources of force fill a line, a surface or a solid region continuously in place of acting from discrete points, a single, double or triple integration takes the place of the summation $\Sigma$. Only the last case occurs in nature itself. Under certain restrictions, however, we can often substitute imagined sources concentrated at points or distributed continuously along lines or over surfaces. We will therefore include those cases in our investigation. We can speak without difficulty of \textit{masses} distributed over a surface or a line, or concentrated at a point, since the term mass here simply means that from which the forces of attraction or repulsion are supposed to act.

\Article{3}
For each point in space, then, we use $x$, $y$, $z$ for its rectangular coordinates and $V$ for the sum of all the mass particles exerting force, \SourceLeaf{206}each divided by its distance from that point. Negative mass particles may be excluded or admitted, depending on the conditions of the particular investigation. Then $V$ is a function of $x$, $y$, $z$, and investigating its characteristic properties will be the key to the theory of attractive or repulsive forces itself. To make the investigations involved easier to handle, we will take the liberty of giving this $V$ a name of its own and calling the quantity the \textit{potential} of the masses to which it refers. This narrower definition suffices for our present investigation. In a broader sense, both for laws of attraction other than the inverse-square law and for the fourth case mentioned in Article 1, the potential could mean the function of $x$, $y$, $z$ whose partial derivatives give the components of the force produced.

For the total force at the point $x$, $y$, $z$ we write $p$, and for the angles its direction makes with the three coordinate axes we write $\alpha$, $\beta$, $\gamma$. The three components are then
\[
p\cos\alpha=\varepsilon\frac{\mathrm{d}V}{\mathrm{d}x},\qquad p\cos\beta=\varepsilon\frac{\mathrm{d}V}{\mathrm{d}y},\qquad p\cos\gamma=\varepsilon\frac{\mathrm{d}V}{\mathrm{d}z}
\]
and
\[
p=\BareRoot{((\frac{\mathrm{d}V}{\mathrm{d}x})^2+(\frac{\mathrm{d}V}{\mathrm{d}y})^2+(\frac{\mathrm{d}V}{\mathrm{d}z})^2)}
\]
\Article{4}
If $\mathrm{d}s$ is the element of any straight or curved line, then $\frac{\mathrm{d}x}{\mathrm{d}s}$, $\frac{\mathrm{d}y}{\mathrm{d}s}$, $\frac{\mathrm{d}z}{\mathrm{d}s}$ are the cosines of the angles this element makes with the coordinate axes. Thus if $\theta$ denotes the angle between the direction of the element and that of the resultant force there, then
\[
\cos\theta=\frac{\mathrm{d}x}{\mathrm{d}s}.\cos\alpha+\frac{\mathrm{d}y}{\mathrm{d}s}.\cos\beta+\frac{\mathrm{d}z}{\mathrm{d}s}.\cos\gamma
\]
The force projected onto the direction of $\mathrm{d}s$ is therefore
\[
p\cos\theta=\varepsilon(\frac{\mathrm{d}V}{\mathrm{d}x}.\frac{\mathrm{d}x}{\mathrm{d}s}+\frac{\mathrm{d}V}{\mathrm{d}y}.\frac{\mathrm{d}y}{\mathrm{d}s}+\frac{\mathrm{d}V}{\mathrm{d}z}.\frac{\mathrm{d}z}{\mathrm{d}s})=\frac{\varepsilon\mathrm{d}V}{\mathrm{d}s}
\]
If we draw a surface through all points at which the potential $V$ has a constant value, it will generally separate the parts of
\SourceLeaf{207}
space where $V$ is less than that value from those where $V$ is greater. If the line $s$ lies on this surface, or at least is tangent to it at the element $\mathrm{d}s$, then $\frac{\mathrm{d}V}{\mathrm{d}s}=0$. Unless the constituents of the total force cancel one another at that point, or in other words $p=0$, in which case the force has no direction, we must have $\cos\theta=0$. From this we conclude that at each point of such a surface the resultant force is normal to the surface itself. It points toward the side of space adjoining the surface where the values of $V$ are greater if $\varepsilon=+1$, and toward the opposite side if $\varepsilon=-1$. We call such a surface a \textit{surface of equilibrium}. Since such a surface can be drawn through every point, the line $s$, unless it lies entirely on a single surface of equilibrium, will meet a different one at each of its points. If $s$ cuts all the surfaces of equilibrium at right angles, a tangent to that line gives the direction of the force everywhere, and $\frac{\mathrm{d}V}{\mathrm{d}s}$ gives its magnitude.

The integral $\int p\cos\theta.\mathrm{d}s$, taken over any segment of the line $s$, is evidently $=\varepsilon(V'-V^0)$, where $V^0$, $V'$ denote the values of the potential at the starting point and the endpoint. Thus if $s$ is a closed line, that integral taken over the entire line is $=0$.

\Article{5}
It is evident that the potential must have an assignable value at every point in space lying \textit{outside} all the attracting or repelling particles. The same is true of its derivatives, both of the first and of higher orders, since on that assumption they too take the form of sums of assignable parts or integrals of differentials whose coefficients all have assignable values. Thus we obtain
\[
\begin{array}{ll}
\displaystyle\frac{\mathrm{d}V}{\mathrm{d}x}=\Sigma\frac{(a-x)\mu}{r^3},&\displaystyle\frac{\mathrm{d}\mathrm{d}V}{\mathrm{d}x^2}=\Sigma(\frac{3(a-x)^2}{r^5}-\frac{1}{r^3})\mu\\[6pt]
\displaystyle\frac{\mathrm{d}V}{\mathrm{d}y}=\Sigma\frac{(b-y)\mu}{r^3},&\displaystyle\frac{\mathrm{d}\mathrm{d}V}{\mathrm{d}y^2}=\Sigma(\frac{3(b-y)^2}{r^5}-\frac{1}{r^3})\mu\\[6pt]
\displaystyle\frac{\mathrm{d}V}{\mathrm{d}z}=\Sigma\frac{(c-z)\mu}{r^3},&\displaystyle\frac{\mathrm{d}\mathrm{d}V}{\mathrm{d}z^2}=\Sigma(\frac{3(c-z)^2}{r^5}-\frac{1}{r^3})\mu
\end{array}
\]
\SourceLeaf{208}
The familiar equation
\[
\frac{\mathrm{d}\mathrm{d}V}{\mathrm{d}x^2}+\frac{\mathrm{d}\mathrm{d}V}{\mathrm{d}y^2}+\frac{\mathrm{d}\mathrm{d}V}{\mathrm{d}z^2}=0
\]
therefore holds for all points in space lying outside the masses exerting the force.

\Article{6}
Of the various cases in which the value of the potential $V$ or its derivatives is sought at a point that does not lie outside the masses exerting the force, we will first consider the case that occurs in nature. In this case the masses fill a definite solid region with uniform or nonuniform density, finite everywhere.

Let $t$ be the entire region containing mass and $\mathrm{d}t$ an infinitesimal element of it, with coordinates $a$, $b$, $c$ and mass element $k\mathrm{d}t$. Let $V$ be the potential at the point $O$, with coordinates $x$, $y$, $z$. The distance from that element to this point is then
\[
\BareRoot{((a-x)^2+(b-y)^2+(c-z)^2)}=r
\]
Consequently we have
\[
V=\int\frac{k\mathrm{d}t}{r}
\]
with the integral taken over the whole region $t$, which entails a triple integration. It is easy to see that the integration can indeed be carried out even if $O$ lies inside the region, although $\frac{1}{r}$ then becomes infinitely large for elements infinitely close to $O$. For if we introduce polar coordinates in place of $a$, $b$, $c$ by setting
\[
a=x+r\cos u,\qquad b=y+r\sin u\cos\lambda,\qquad c=z+r\sin u\sin\lambda
\]
we obtain $\mathrm{d}t=rr\sin u.\mathrm{d}u.\mathrm{d}\lambda.\mathrm{d}r$ and hence
\[
V=\int\!\int\!\int kr\sin u.\mathrm{d}u.\mathrm{d}\lambda.\mathrm{d}r
\]
where the integration with respect to $r$ must run from $r=0$ to the value at the boundary of $t$, from $\lambda=0$ to $\lambda=2\pi$, and from $u=0$ to $u=\pi$. Thus $V$ must have a definite finite value.

\SourceLeaf{209}
It is also easy to see that here too we may set
\[
\frac{\mathrm{d}V}{\mathrm{d}x}=\int k\mathrm{d}t.\frac{\mathrm{d}\frac{1}{r}}{\mathrm{d}x}=\int\frac{k(a-x)\mathrm{d}t}{r^3}=X
\]
This is justified because this expression too, which becomes
\[
\int\!\int\!\int k\cos u.\sin u.\mathrm{d}u.\mathrm{d}\lambda.\mathrm{d}r
\]
when polar coordinates are used, can indeed be integrated. Thus $X$ has a definite finite value that varies continuously, since all the elements infinitely close to $O$ make only an infinitesimal contribution to it. For similar reasons we may also set
\[
\begin{array}{l}
\displaystyle\frac{\mathrm{d}V}{\mathrm{d}y}=\int\frac{k(b-y)\mathrm{d}t}{r^3}=Y\\[6pt]
\displaystyle\frac{\mathrm{d}V}{\mathrm{d}z}=\int\frac{k(c-z)\mathrm{d}t}{r^3}=Z
\end{array}
\]
and these quantities, like $V$, therefore have definite values within $t$ that vary continuously. This remains true on the boundary of $t$.

\Article{7}
As for derivatives of higher orders, a different procedure is needed for points within $t$, since, for example, it is not permissible to transform $\frac{\mathrm{d}X}{\mathrm{d}x}$ into
\[
\int k\mathrm{d}t.\frac{\mathrm{d}\frac{a-x}{r^3}}{\mathrm{d}x}
\]
that is, into
\[
\int k(\frac{3(a-x)^2-rr}{r^5})\mathrm{d}t
\]
since on closer examination this expression would be merely a sign without a definite, clear meaning. In fact, in any part of $t$ that encloses the point, however small, one can find subregions over which this integral exceeds any prescribed value, positive or negative. Thus the essential
\SourceLeaf{210}
condition under which alone the whole integral can be given a clear meaning is missing here, namely the applicability of the method of exhaustion.

\Article{8}
Before undertaking this investigation in its general form, it will help to clarify our ideas by considering a very simple special case.

Let $t$ be a sphere with radius $=R$ and its center at the origin of coordinates. Let the density of the mass filling the sphere be constant $=k$, and let the distance of the point $O$ from the center be denoted by $\rho=\BareRoot{(xx+yy+zz)}$. As is known, the potential has two different expressions, according as $O$ lies inside or outside the sphere. In the first case we have
\[
V=2\pi kRR-\tfrac{2}{3}\pi k\rho\rho=2\pi kRR-\tfrac{2}{3}\pi k(xx+yy+zz)
\]
and in the second
\[
V=\frac{4\pi kR^3}{3\rho}
\]
On the surface of the sphere both expressions give the same value $\tfrac{4}{3}\pi kRR$, and the potential therefore varies continuously throughout space.

For the derivatives we obtain, in the interior,
\[
\begin{array}{l}
\displaystyle\frac{\mathrm{d}V}{\mathrm{d}x}=X=-\tfrac{4}{3}\pi kx\\[6pt]
\displaystyle\frac{\mathrm{d}V}{\mathrm{d}y}=Y=-\tfrac{4}{3}\pi ky\\[6pt]
\displaystyle\frac{\mathrm{d}V}{\mathrm{d}z}=Z=-\tfrac{4}{3}\pi kz
\end{array}
\]
and in the exterior,
\[
\begin{array}{l}
\displaystyle X=-\frac{4\pi kR^3x}{3\rho^3}\\[6pt]
\displaystyle Y=-\frac{4\pi kR^3y}{3\rho^3}\\[6pt]
\displaystyle Z=-\frac{4\pi kR^3z}{3\rho^3}
\end{array}
\]
Here too the latter formulas give the same values on the surface as the former, so that $X$, $Y$, $Z$ also vary continuously throughout space.

\SourceLeaf{211}
The derivatives of these quantities, however, behave differently. In the interior we have
\[
\frac{\mathrm{d}X}{\mathrm{d}x}=-\tfrac{4}{3}\pi k,\qquad \frac{\mathrm{d}Y}{\mathrm{d}y}=-\tfrac{4}{3}\pi k,\qquad \frac{\mathrm{d}Z}{\mathrm{d}z}=-\tfrac{4}{3}\pi k
\]
and in the exterior,
\[
\begin{array}{l}
\displaystyle\frac{\mathrm{d}X}{\mathrm{d}x}=\frac{4\pi kR^3(3xx-\rho\rho)}{3\rho^5}\\[6pt]
\displaystyle\frac{\mathrm{d}Y}{\mathrm{d}y}=\frac{4\pi kR^3(3yy-\rho\rho)}{3\rho^5}\\[6pt]
\displaystyle\frac{\mathrm{d}Z}{\mathrm{d}z}=\frac{4\pi kR^3(3zz-\rho\rho)}{3\rho^5}
\end{array}
\]
On the surface the latter values exceed the former, respectively, by
\[
\frac{4\pi kxx}{RR},\qquad \frac{4\pi kyy}{RR},\qquad \frac{4\pi kzz}{RR}
\]
Thus these derivatives vary continuously throughout the interior and throughout the exterior, but jump on passing from one to the other. On the separating surface itself, they must be assigned two values, according as $\mathrm{d}x$, $\mathrm{d}y$, $\mathrm{d}z$ are regarded as positive or negative.

The same happens with the other six derivatives
\[
\frac{\mathrm{d}X}{\mathrm{d}y},\qquad \frac{\mathrm{d}X}{\mathrm{d}z},\qquad \frac{\mathrm{d}Y}{\mathrm{d}x},\qquad \frac{\mathrm{d}Y}{\mathrm{d}z},\qquad \frac{\mathrm{d}Z}{\mathrm{d}x},\qquad \frac{\mathrm{d}Z}{\mathrm{d}y}
\]
which are all $=0$ inside the sphere and undergo the jumps
\[
\frac{4\pi kxy}{RR},\qquad \frac{4\pi kxz}{RR}\text{ etc.}
\]
on crossing the spherical surface.

The sum $\frac{\mathrm{d}X}{\mathrm{d}x}+\frac{\mathrm{d}Y}{\mathrm{d}y}+\frac{\mathrm{d}Z}{\mathrm{d}z}$ or $\frac{\mathrm{d}\mathrm{d}V}{\mathrm{d}x^2}+\frac{\mathrm{d}\mathrm{d}V}{\mathrm{d}y^2}+\frac{\mathrm{d}\mathrm{d}V}{\mathrm{d}z^2}$ is $=-4\pi k$ inside the sphere and $=0$ in the exterior. On the surface itself, however, it loses its simple meaning. Strictly speaking, we can only say that it is a sum of three parts, each with two different values. Thus there are actually eight combinations. One agrees with the value on the inner side and another with the value on the outer side, while the other six have no meaning at all. I cannot admit the validity of the ana%
\SourceLeaf{212}lysis by which some geometers have obtained the value $-2\pi k$ on the surface of the sphere, or the mean of the interior and exterior values, insofar as the concept of a differential quotient is understood in its mathematical purity.

\Article{9}
The result found in the preceding example is only a particular case of the general theorem that, when the point $O$ lies inside the mass exerting the force, the value of $\frac{\mathrm{d}\mathrm{d}V}{\mathrm{d}x^2}+\frac{\mathrm{d}\mathrm{d}V}{\mathrm{d}y^2}+\frac{\mathrm{d}\mathrm{d}V}{\mathrm{d}z^2}$ equals the product of $-4\pi$ and the density at $O$. The following seems the most satisfactory way to establish this important theorem.

We assume that the density $k$ has no jumps anywhere within $t$, or that it is a function $f(a,b,c)$ of $a$, $b$, $c$ whose value varies continuously everywhere inside $t$ and outside $t$ is $=0$.

Let $t'$ be the region into which $t$ passes when the first coordinate of every point of the boundary surface is diminished by $e$, or, equivalently, when the boundary surface is moved backward by $e$ parallel to the first coordinate axis. Let $t$ consist of the regions $t^0$ and $\theta$, and $t'$ of $t^0$ and $\theta'$, so that $t^0$ is the whole region common to $t$ and $t'$. We consider the three integrals
\[
\begin{array}{ll}
\displaystyle\int\frac{f(a,b,c)(a-x)\mathrm{d}t}{((a-x)^2+(b-y)^2+(c-z)^2)^{\frac{3}{2}}}&.\ .\ .\ .\ .\ .\ .\ .\quad(1)\\[8pt]
\displaystyle\int\frac{f(a,b,c)(a-x-e)\mathrm{d}t}{((a-x-e)^2+(b-y)^2+(c-z)^2)^{\frac{3}{2}}}&.\ .\ .\ .\ .\ .\ .\quad(2)\\[8pt]
\displaystyle\int\frac{f(a+e,b,c)(a-x)\mathrm{d}t}{((a-x)^2+(b-y)^2+(c-z)^2)^{\frac{3}{2}}}&.\ .\ .\ .\ .\ .\ .\ .\quad(3)
\end{array}
\]
where integral (1), taken over the whole region $t$, gives the value of $\frac{\mathrm{d}V}{\mathrm{d}x}$ or $X$ at the point $O$. Integral (2), also taken over all of $t$, gives the value of $\frac{\mathrm{d}V}{\mathrm{d}x}$ at the point with coordinates $x+e$, $y$, $z$, a value we will denote by $X+\xi$. Evidently this integral is entirely identical to integral (3) taken over the whole region $t'$. Thus if
\SourceLeaf{213}%
\[
\begin{array}{rlcl}
\text{integral (1), taken over}&t^0&.\ .\ .\ .\ .\ .\ .\ .&l\\
\text{over}&\theta&.\ .\ .\ .\ .\ .\ .\ .&\lambda\\
\text{integral (3), taken over}&t^0&.\ .\ .\ .\ .\ .\ .\ .&l'\\
\text{over}&\theta'&.\ .\ .\ .\ .\ .\ .\ .&\lambda'
\end{array}
\]
we have $X=l+\lambda$, $X+\xi=l'+\lambda'$.

If we set $f(a+e,b,c)-f(a,b,c)=\Delta k$, the integral
\[
\int\frac{\frac{\Delta k}{e}(a-x)\mathrm{d}t}{((a-x)^2+(b-y)^2+(c-z)^2)^{\frac{3}{2}}}\quad .\ .\ .\ .\ .\ .\ .\ .\quad(4)
\]
taken over $t^0$ is $=\frac{l'-l}{e}$.

The results so far hold generally for every position of $O$. In the further argument, the case where $O$ lies on the surface itself is to be excluded, or we are to assume that $O$ lies at a measurable distance from the surface, inside or outside $t$.

Now let $e$ become infinitely small. The regions $\theta$, $\theta'$ are then two infinitely thin layers adjoining the surface of $t$. Divide this surface into elements $\mathrm{d}s$, and use $\alpha$ for the angle an outward normal erected at $\mathrm{d}s$ makes with the first coordinate axis. Then $\alpha$ is evidently acute wherever the surface of $t$ adjoins $\theta$ and obtuse wherever it adjoins $\theta'$. The elements of $\theta$ are therefore expressed by $e\cos\alpha\mathrm{d}s$, and those of $\theta'$ by $-e\cos\alpha\mathrm{d}s$. From this it is easy to conclude that $\frac{\lambda-\lambda'}{e}$ passes into the integral
\[
\int\frac{f(a,b,c)(a-x)\cos\alpha.\mathrm{d}s}{((a-x)^2+(b-y)^2+(c-z)^2)^{\frac{3}{2}}}
\]
or, equivalently,
\[
\int\frac{k(a-x)\cos\alpha.\mathrm{d}s}{r^3}
\]
taken over the entire surface, where $k$ means the density at the element $\mathrm{d}s$.

Assuming an infinitely small value of $e$, the value of $\frac{\Delta k}{e}$ also passes into the value of the partial derivative $\frac{\mathrm{d}f(a,b,c)}{\mathrm{d}a}$ or $\frac{\mathrm{d}k}{\mathrm{d}a}$, and the value of integral (4) or $\frac{(l'-l)}{e}$ into the integral
\SourceLeaf{214}%
\[
\int\frac{\frac{\mathrm{d}k}{\mathrm{d}a}.(a-x)\mathrm{d}t}{r^3}
\]
taken over the entire region $t$.

Finally, for an infinitely small $e$, the expression $\frac{l'-l}{e}-\frac{\lambda-\lambda'}{e}$ or $\frac{\xi}{e}$ is just the value of the partial derivative $\frac{\mathrm{d}X}{\mathrm{d}x}$ or $\frac{\mathrm{d}\mathrm{d}V}{\mathrm{d}x^2}$. We therefore have the simple result
\[
\frac{\mathrm{d}\mathrm{d}V}{\mathrm{d}x^2}=\frac{\mathrm{d}X}{\mathrm{d}x}=\int\frac{\frac{\mathrm{d}k}{\mathrm{d}a}.(a-x)\mathrm{d}t}{r^3}-\int\frac{k(a-x)\cos\alpha.\mathrm{d}s}{r^3}
\]
where the first integration is over the entire region $t$ and the second over its entire surface.

This result holds however close $O$ lies to the surface on the inside or the outside, provided the point lies off the surface itself, where $\frac{\mathrm{d}X}{\mathrm{d}x}$ will have two different values. The first integral varies continuously on crossing the surface. According to a theorem to be proved below, however, $-\int\frac{k(a-x)\cos\alpha\mathrm{d}s}{r^3}$ changes by the finite amount $4\pi k\cos\alpha$ on passing from an interior point infinitely close to the surface to an exterior point, where $k$ and $\alpha$ refer to the crossing point. The difference between the two values of $\frac{\mathrm{d}X}{\mathrm{d}x}$ there will be just as great.

\Article{10}
In the same way, if $\beta$ and $\gamma$ have the same meaning with respect to the second and third coordinate axes as $\alpha$ has with respect to the first, and the position of $O$ is subject to the same restriction as before, we obtain
\[
\begin{aligned}
\frac{\mathrm{d}Y}{\mathrm{d}y}&=\int\frac{\frac{\mathrm{d}k}{\mathrm{d}b}(b-y)\mathrm{d}t}{r^3}-\int\frac{k(b-y)\cos\beta.\mathrm{d}s}{r^3}\\
\frac{\mathrm{d}Z}{\mathrm{d}z}&=\int\frac{\frac{\mathrm{d}k}{\mathrm{d}c}(c-z)\mathrm{d}t}{r^3}-\int\frac{k(c-z)\cos\gamma.\mathrm{d}s}{r^3}
\end{aligned}
\]
Now consider that
\[
\frac{\mathrm{d}k}{\mathrm{d}a}.\frac{a-x}{r}+\frac{\mathrm{d}k}{\mathrm{d}b}.\frac{b-y}{r}+\frac{\mathrm{d}k}{\mathrm{d}c}.\frac{c-z}{r}
\]
\SourceLeaf{215}%
is just the value of the derivative $\frac{\mathrm{d}k}{\mathrm{d}r}$, provided that in this differentiation only the length of $r$ is treated as variable and its direction as constant. Also,
\[
\frac{a-x}{r}.\cos\alpha+\frac{b-y}{r}.\cos\beta+\frac{c-z}{r}.\cos\gamma=\cos\psi
\]
if $\psi$ denotes the angle between the outward normal at $\mathrm{d}s$ and the extended straight line $r$. It follows that, if the integral
\[
\int\frac{\frac{\mathrm{d}k}{\mathrm{d}r}}{rr}.\mathrm{d}t
\]
taken over the entire region $t$ is denoted by $M$, and the integral
\[
\int\frac{k\cos\psi}{rr}.\mathrm{d}s
\]
taken over the entire surface of $t$ by $N$, then
\[
\frac{\mathrm{d}\mathrm{d}V}{\mathrm{d}x^2}+\frac{\mathrm{d}\mathrm{d}V}{\mathrm{d}y^2}+\frac{\mathrm{d}\mathrm{d}V}{\mathrm{d}z^2}=M-N
\]

To carry out the first integration, we draw a spherical surface with center $O$ and radius $1$ and divide it into elements $\mathrm{d}\sigma$. The straight lines drawn from $O$ through all points of the perimeter of $\mathrm{d}\sigma$ and extended indefinitely form a conical surface (in the broader sense of the word). This cuts out a region from the whole of $t$, possibly consisting of several separate pieces, with general element $rr\mathrm{d}\sigma.\mathrm{d}r$. The part of $M$ belonging to this region is therefore expressed by $\mathrm{d}\sigma.\int\frac{\mathrm{d}k}{\mathrm{d}r}.\mathrm{d}r$, when the integration is taken over all parts within $t$ of a straight line through $O$ and a point of $\mathrm{d}\sigma$. This line $r$ is extended as far as necessary. Suppose that this straight line cuts the surface of $t$ successively at $O'$, $O''$, $O'''$, $O''''$ and so on. Denote by $r'$, $r''$, $r'''$, $r''''$ and so on the values of $r$ at those points, by $\mathrm{d}s'$, $\mathrm{d}s''$, $\mathrm{d}s'''$, $\mathrm{d}s''''$ and so on the corresponding elements cut from the surface of $t$ by the elementary cone, by $k'$, $k''$, $k'''$, $k''''$ and so on the values of $k$, and by $\psi'$, $\psi''$, $\psi'''$, $\psi''''$ and so on the values of $\psi$ at those elements. It is then easy to see that
\SourceLeaf{216}%

I. In the case where $O$ lies inside $t$, the number of those points is odd, and the integration $\int\frac{\mathrm{d}k}{\mathrm{d}r}.\mathrm{d}r$ must be carried out from $r=0$ to $r=r'$, then from $r=r''$ to $r=r'''$ and so on. Thus if the density at $O$ is denoted by $k^0$, we obtain
\[
\int\frac{\mathrm{d}k}{\mathrm{d}r}.\mathrm{d}r=-k^0+k'-k''+k'''-k''''+\text{ etc.}
\]
Since the angles $\psi'$, $\psi''$, $\psi'''$, $\psi''''$ and so on are evidently alternately acute and obtuse, we have
\[
\begin{aligned}
\mathrm{d}s'.\cos\psi'&=+r'r'\mathrm{d}\sigma\\
\mathrm{d}s''.\cos\psi''&=-r''r''\mathrm{d}\sigma\\
\mathrm{d}s'''.\cos\psi'''&=+r'''r'''\mathrm{d}\sigma\\
\mathrm{d}s''''.\cos\psi''''&=-r''''r''''\mathrm{d}\sigma
\end{aligned}
\]
and so on, and consequently
\[
\begin{aligned}
\mathrm{d}\sigma\int\frac{\mathrm{d}k}{\mathrm{d}r}.\mathrm{d}r&=-k^0\mathrm{d}\sigma+\frac{k'\cos\psi'}{r'r'}\mathrm{d}s'+\frac{k''\cos\psi''}{r''r''}\mathrm{d}s''+\frac{k'''\cos\psi'''}{r'''r'''}+\text{ etc.}\\
&=-k^0\mathrm{d}\sigma+\sum\frac{k\cos\psi}{rr}\mathrm{d}s
\end{aligned}
\]
where the sum extends over all $\mathrm{d}s$ corresponding to the element $\mathrm{d}\sigma$. Integrating over all $\mathrm{d}\sigma$ therefore gives
\[
M=-4\pi k^0+\int\frac{k\cos\psi}{rr}\mathrm{d}s
\]
where the integral must be taken over the entire surface, or $M=-4\pi k^0+N$. Consequently
\[
\frac{\mathrm{d}\mathrm{d}V}{\mathrm{d}x^2}+\frac{\mathrm{d}\mathrm{d}V}{\mathrm{d}y^2}+\frac{\mathrm{d}\mathrm{d}V}{\mathrm{d}z^2}=-4\pi k^0
\]

II. In the case where $O$ lies outside $t$, only those $\mathrm{d}\sigma$ need be considered for which the straight line through $O$ and a point of $\mathrm{d}\sigma$ actually meets the region $t$. Here the number of points $O'$, $O''$, $O'''$ and so on is always even, and the angles $\psi'$, $\psi''$, $\psi'''$ and so on are alternately obtuse and acute. Thus
\[
\begin{aligned}
\mathrm{d}s'.\cos\psi'&=-r'r'\mathrm{d}\sigma,\\
\mathrm{d}s''.\cos\psi''&=+r''r''\mathrm{d}\sigma,\\
\mathrm{d}s'''\cos\psi'''&=-r'''r'''\mathrm{d}\sigma\text{ etc.}
\end{aligned}
\]
Since the integration $\int\frac{\mathrm{d}k}{\mathrm{d}r}.\mathrm{d}r$ here must be carried out from $r=r'$ to $r=r''$, then from $r=r'''$
\SourceLeaf{217}%
to $r=r''''$ and so on, we obtain
\[
\begin{aligned}
\mathrm{d}\sigma.\int\frac{\mathrm{d}k}{\mathrm{d}r}.\mathrm{d}r&=\frac{k'\cos\psi'}{r'r'}.\mathrm{d}s'+\frac{k''\cos\psi''}{r''r''}.\mathrm{d}s''\\
&\quad{}+\frac{k'''\cos\psi'''}{r'''r'''}.\mathrm{d}s'''+\text{etc.}\\
&=\sum\frac{k\cos\psi}{rr}\mathrm{d}s
\end{aligned}
\]
and after the second integration over all the relevant $\mathrm{d}\sigma$,
\[
M=\int\frac{k\cos\psi}{rr}\mathrm{d}s=N
\]
hence, as is already known,
\[
\frac{\mathrm{d}\mathrm{d}V}{\mathrm{d}x^2}+\frac{\mathrm{d}\mathrm{d}V}{\mathrm{d}y^2}+\frac{\mathrm{d}\mathrm{d}V}{\mathrm{d}z^2}=0
\]

\Article{11}
Although our proof assumes that the density varies continuously throughout the \textit{whole} region $t$, that condition is not necessary for our result to hold. It is enough for the density at $O$ to vary continuously in every direction, or for $O$ to lie inside a region, however small, satisfying this condition. If we set the potential of the mass contained in \textit{this} region $=V'$ and the potential of the remaining masses outside it $=V''$, the total potential is $V=V'+V''$. Since, by the preceding article,
\[
\begin{aligned}
\frac{\mathrm{d}\mathrm{d}V'}{\mathrm{d}x^2}+\frac{\mathrm{d}\mathrm{d}V'}{\mathrm{d}y^2}+\frac{\mathrm{d}\mathrm{d}V'}{\mathrm{d}z^2}&=-4\pi k^0\\
\frac{\mathrm{d}\mathrm{d}V''}{\mathrm{d}x^2}+\frac{\mathrm{d}\mathrm{d}V''}{\mathrm{d}y^2}+\frac{\mathrm{d}\mathrm{d}V''}{\mathrm{d}z^2}&=0
\end{aligned}
\]
we have
\[
\frac{\mathrm{d}\mathrm{d}V}{\mathrm{d}x^2}+\frac{\mathrm{d}\mathrm{d}V}{\mathrm{d}y^2}+\frac{\mathrm{d}\mathrm{d}V}{\mathrm{d}z^2}=-4\pi k^0
\]
If this condition fails at the point $O$, so that it lies on the surface separating two regions in which the density varies continuously within each region but jumps on passing from one to the other, then, generally speaking, $\frac{\mathrm{d}\mathrm{d}V}{\mathrm{d}x^2}$, $\frac{\mathrm{d}\mathrm{d}V}{\mathrm{d}y^2}$, $\frac{\mathrm{d}\mathrm{d}V}{\mathrm{d}z^2}$ each have two different values there. Their sum is subject to the same observation made at the end of Article 8.
\SourceLeaf{218}%
\Article{12}
As already noted above, we also include in our investigation the ideal case in which forces of attraction or repulsion are assumed to come from the parts of a \textit{surface}. We take the liberty of describing the source of the force as a mass distributed over the surface. By the density at any point of the surface we mean, in this case, the quotient obtained by dividing the mass contained in a surface element to which the point belongs by that element. This density may be uniform, the same at all points, or nonuniform. In the latter case it may vary continuously over the whole surface, that is, differ only infinitely little at any two infinitely close points. Or the whole surface may break up into two or more pieces, with continuous variation within each piece and a jump on passing from one to another. We can also imagine a distribution in which the density becomes infinitely great at particular points or lines, while the total mass remains finite. The surface itself, insofar as it is not a plane, will generally be assigned continuous curvature, while allowing interruptions at particular points (corners) or lines (edges).

\expandafter\def\csname HistoricalNoteText22\endcsname{It is easy to establish that the integral expressing the potential is finite by dividing the surface into elements in a manner similar to that used in Article 15 below. At the same time this shows that the parts of the surface infinitely close to the two points in question contribute only infinitely little to the whole integral, from which the preceding assertion is easily proved.}%
Under these assumptions the potential also has a definite finite value at each point of the surface itself where the density is not infinitely great. The value at a second point infinitely close to it, whether on the surface or off it, can differ from that value only infinitely little\HistoricalNote{22}. In other words, the potential varies continuously along any line, whether it lies on the surface itself or crosses it.

\Article{13}
If we denote by $k$ the density in the surface element $\mathrm{d}s$, by
\SourceLeaf{219}%
$a$, $b$, $c$ the coordinates of a point belonging to it, by $r$ its distance from a point $O$ whose coordinates are $x$, $y$, $z$, and by $V$ the potential of the mass contained in the surface at the point $O$, then $V=\int\frac{k\mathrm{d}s}{r}$, with the integral taken over the entire surface. Finally, if we denote by $X$, $Y$, $Z$ the integrals
\[
\int\frac{k(a-x)\mathrm{d}s}{r^3},\quad\int\frac{k(b-y)\mathrm{d}s}{r^3},\quad\int\frac{k(c-z)\mathrm{d}s}{r^3}
\]
understood in the same way, then $X$, $Y$, $Z$ do have exactly the same meaning as $\frac{\mathrm{d}V}{\mathrm{d}x}$, $\frac{\mathrm{d}V}{\mathrm{d}y}$, $\frac{\mathrm{d}V}{\mathrm{d}z}$ as long as $O$ lies off the surface. Strictly speaking, however, this is no longer true when $O$ is a point of the surface itself, and the discrepancy takes different forms according to the angle that the normal to the surface makes with the coordinate axis in question. It is evidently enough to describe the behavior with respect to the first coordinate axis alone.

I. If that angle $=0$, then at $O$ the integral $X$ has a definite value, whereas $\frac{\mathrm{d}V}{\mathrm{d}x}$ has two different values according to whether $\mathrm{d}x$ is regarded as positive or negative.

II. If the angle is a right angle, the expression for $X$ cannot be integrated in the true sense (since a remark similar to that in Article 7 then applies), whereas $\frac{\mathrm{d}V}{\mathrm{d}x}$ has only one definite value.

III. If the angle is acute, $X$ behaves as in the second case, and $\frac{\mathrm{d}V}{\mathrm{d}x}$ as in the first.

Further special modifications occur if continuity is interrupted at $O$ with respect to either the density or the curvature. For our main purpose, however, it is unnecessary to discuss in detail such exceptional cases, which can occur only along individual lines or at individual points. In examining the matter more closely, we shall therefore assume that there is a definite finite density and a definite tangent plane at the point in question.

\Article{14}
Before undertaking the investigation in its generality, it will be useful to consider a simple special case. Let the surface be a portion $A$ of a spherical surface, and let its density be uniform, or $k$ constant. Thus $V$, $X$ are the values of the integrals
\SourceLeaf{220}%
\[
\int\frac{k\mathrm{d}s}{r},\quad\int\frac{k(a-x)\mathrm{d}s}{r^3}
\]
taken over $A$. If we denote by $V'$, $X'$ the same integrals taken over the remaining part of the spherical surface $B$, and by $V^0$, $X^0$ those taken over the entire spherical surface, then $V=V^0-V'$, $X=X^0-X'$. We shall also denote the radius of the sphere by $R$, place the origin of coordinates at the center of the sphere, and set $\BareRoot{(xx+yy+zz)}$, or the distance of the point $O$ from the center of the sphere, $=\rho$.

It is known that $V^0=4\pi kR$ when $O$ lies inside the sphere, whereas $V^0=\frac{4\pi kRR}{\rho}$ when $O$ lies outside it. On the spherical surface itself the two values coincide. The derivative $\frac{\mathrm{d}V^0}{\mathrm{d}x}$ is therefore $=0$ inside the sphere and $=-\frac{4\pi kRRx}{\rho^3}$ outside it. On the spherical surface itself, however, both values apply at once, according to the sign of $\mathrm{d}x$. These two values are equal only when $x=0$, which corresponds to case II of the preceding article.

The expression for $X^0$, which has the same meaning as $\frac{\mathrm{d}V^0}{\mathrm{d}x}$ inside and outside the sphere, becomes an empty symbol on the surface, insofar as a true integration is not permissible. The sole exception is the case in which, for the infinitely close elements of the surface, $a-x$ becomes an infinitesimal of a higher order than $r$, namely when $y=0$, $z=0$. $x=\pm R$, in which case integration gives $X^0=\mp2\pi k$, agreeing with neither value of $\frac{\mathrm{d}V^0}{\mathrm{d}x}$, but rather with their mean. This case evidently belongs to I in the preceding article.

Now, if $O$ is a point on the spherical surface lying inside $A$, then $X'$ and $\frac{\mathrm{d}V'}{\mathrm{d}x}$ have the same meaning and have definite values that vary continuously. It follows that the relation between $X^0-X'$ and $\frac{\mathrm{d}V^0}{\mathrm{d}x}-\frac{\mathrm{d}V'}{\mathrm{d}x}$, that is, between $X$ and $\frac{\mathrm{d}V}{\mathrm{d}x}$, is exactly the same as that between $X^0$ and $\frac{\mathrm{d}V^0}{\mathrm{d}x}$. The propositions stated in the preceding article therefore follow of themselves.

\Article{15}
For the more general investigation, it is advantageous to place the origin of coordinates at a point $P$ on the surface itself and to make the first coordinate axis perpendicular to the tangent plane at $P$.
\SourceLeaf{221}%

If we denote by $\psi$ the angle between the normal to an arbitrary surface element $\mathrm{d}s$ and the first coordinate axis, then $\cos\psi.\mathrm{d}s$ is the projection of $\mathrm{d}s$ onto the plane of $b$ and $c$. If we set
\[
\BareRoot{(bb+cc)}=\rho,\quad b=\rho\cos\theta,\quad c=\rho\sin\theta
\]
then $\rho\mathrm{d}\rho.\mathrm{d}\theta$ represents an arbitrary element of this plane, and the corresponding surface element is $\mathrm{d}s=\frac{\rho\mathrm{d}\rho.\mathrm{d}\theta}{\cos\psi}$. The mass element contained in it is therefore $=h\rho\mathrm{d}\rho.\mathrm{d}\theta$, if for brevity we write $h$ for $\frac{k}{\cos\psi}$.

We shall now investigate how far the value of $X$ changes by a jump as the point $O$ passes along the first coordinate axis from one side of the surface to the other, or as $x$ passes from a negative value to a positive one. For this question it evidently makes no difference whether we consider the entire surface or only an arbitrarily small portion enclosing the point $P$, since the contribution of the remaining portion of the surface to the value of $X$ varies continuously. We may therefore take $\rho$ only from $0$ to an arbitrarily small limiting value $\rho'$ and assume that, within the surface so bounded, $h$ and $\frac{a}{\rho}$ vary continuously everywhere. For each definite value of $\theta$, if we set the value of the integral $\int\frac{h(a-x)\rho\mathrm{d}\rho}{r^3}$, taken from $\rho=0$ to $\rho=\rho'$, $=Q$, then $X=\int Q\mathrm{d}\theta$, where the integration is to extend from $\theta=0$ to $\theta=2\pi$.

We must now compare the values of $X$ for $x=0$, for an infinitely small positive $x$, and for an infinitely small negative one (the other two coordinates $y$, $z$ being taken $=0$ throughout). We denote these three values of $X$ by $X^0$, $X'$, $X''$, and the corresponding values of $Q$ by $Q^0$, $Q'$, $Q''$.

Since $r=\BareRoot{((a-x)^2+\rho\rho)}$, treating $\theta$ as constant gives
\[
\mathrm{d}\frac{h(a-x)}{r}=-\frac{h(a-x)\rho\mathrm{d}\rho}{r^3}+\frac{\mathrm{d}h}{\mathrm{d}\rho}.\frac{a-x}{r}.\mathrm{d}\rho+\frac{\mathrm{d}a}{\mathrm{d}\rho}.\frac{h\rho\rho}{r^3}.\mathrm{d}\rho
\]
and consequently
\[
Q=\int\frac{\mathrm{d}h}{\mathrm{d}\rho}.\frac{a-x}{r}.\mathrm{d}\rho+\int\frac{\mathrm{d}a}{\mathrm{d}\rho}.\frac{h\rho\rho}{r^3}.\mathrm{d}\rho-\frac{h'(a'-x)}{r'}+\text{Const.}
\]
where both integrations extend from $\rho=0$ to $\rho=\rho'$, and the values
\SourceLeaf{222}%
of $h,a,r$ for $\rho=\rho'$ are denoted by $h',a',r'$. As the constant we must take the value of $\frac{h(a-x)}{r}$ for $\rho=0$. If we denote the density at $P$ by $k^0$, this value is $=-k^0$ for a positive $x$ and $=+k^0$ for a negative one, since for $\rho=0$ we evidently have $a=0$, $\psi=0$, $h=k^0$, $x=\pm r$. For the case $x=0$, however, the constant must be the limiting value of $\frac{ha}{r}$ as $\rho$ becomes infinitely small. This is $=0$, because $a$ becomes an infinitesimal of a higher order than $r$.

The value of the integral $\int\frac{\mathrm dh}{\mathrm d\rho}.\frac{a-x}{r}.\mathrm d\rho$ remains the same up to an infinitely small difference, whether we set $x=0$ or make it infinitely small $=\pm\varepsilon$. For if we split that integral into
\[
\int_0^\delta\frac{\mathrm dh}{\mathrm d\rho}.\frac{a-x}{r}.\mathrm d\rho+\int_\delta^{\rho'}\frac{\mathrm dh}{\mathrm d\rho}.\frac{a-x}{r}.\mathrm d\rho
\]
it is clear that the assertion holds for the first part if $\delta$ is infinitely small, and for the second if $\frac{\delta}{\varepsilon}$ is infinitely large, hence for the whole if $\delta$ is an infinitesimal of a lower order than $\varepsilon$.

A similar argument also applies to the integral $\int\frac{\mathrm da}{\mathrm d\rho}.\frac{h\rho\rho}{r^3}.\mathrm d\rho$, if the points of the surface corresponding to a definite value of $\theta$ form a curve with measurable curvature at $P$, so that $\frac{a}{\rho\rho}$ has a finite, continuously varying value in the region considered here. If we denote this value by $A$, then
\[
\frac{\mathrm da}{\mathrm d\rho}=2A\rho+\frac{\mathrm dA}{\mathrm d\rho}.\rho\rho
\]
and that integral splits into the following two
\[
\int\frac{2\rho^3Ah\mathrm d\rho}{r^3}+\int\frac{\mathrm dA}{\mathrm d\rho}.\frac{\rho^4}{r^3}h\mathrm d\rho
\]
for both of which the validity of the preceding argument is evident.

Finally, the values of $\frac{h'(a'-x)}{r'}$ are also evidently equal, up to infinitely small differences, for all three values of $x$.

It follows that $Q'+k^0$, $Q^0$, $Q''-k^0$ are equal up to infinitely small differences, and the same must therefore hold for
\[
\int(Q'+k^0)\mathrm d\theta,\quad\int Q^0\mathrm d\theta,\quad\int(Q''-k^0)\mathrm d\theta
\]
or for the quantities
\[
X'+2\pi k^0,\quad X^0,\quad X''-2\pi k^0
\]

\SourceLeaf{223}%
This important proposition can also be expressed as follows. The limiting value of $X$ as positive $x$ becomes infinitely small is $X^0-2\pi k^0$, whereas as negative $x$ becomes infinitely small it is $X^0+2\pi k^0$. In other words, $X$ changes twice by a jump of $-2\pi k^0$ as $x$ passes from a negative value to a positive one, the first time when $x$ reaches the value $0$, and the second when it passes beyond it.

\Article{16}
The proof in the preceding article assumes that the intersections of the surface with planes drawn through the first coordinate axis have measurable curvature at $P$. Our result remains valid even if the curvature at $P$ is infinitely large, with a single exception. That $\frac a\rho$ must itself become infinitely small for infinitely small $\rho$ follows already from the assumption that the surface has a definite tangent plane at $P$. The two quantities are of the same order only when the radius of curvature is finite. For an infinitely small radius of curvature, however, $\frac a\rho$ will be of lower order than $\rho$. We will now show that our results remain valid in the latter case too, provided the orders of the two quantities are \textit{comparable}.

Suppose, then, that $\frac a\rho$ is of the same order as $\rho^\mu$, where $\mu$ is a finite positive exponent. Thus $\frac{a}{\rho^{1+\mu}}$ is a finite quantity that varies continuously in the region under consideration. We will denote this quantity by $B$. The integral $\int\frac{\mathrm da}{\mathrm d\rho}.\frac{h\rho\rho}{r^3}\mathrm d\rho$ therefore splits into the following two
\[
\int\frac{(1+\mu)\rho^{2+\mu}hB\mathrm d\rho}{r^3}+\int\frac{\rho^{3+\mu}}{r^3}.\frac{\mathrm dB}{\mathrm d\rho}.h\mathrm d\rho
\]
The conclusions of the preceding article apply directly to the second integral and to the first after a slight transformation. If we set $\frac1\mu=m$, $\rho^\mu=\sigma$ or $\rho=\sigma^m$, that integral becomes
\[
=(m+1)\int\frac{Bh\sigma^{3m}\mathrm d\sigma}{(\sigma^{2m}+(a-x)^2)^{\frac32}}
\]
This integral, too, evidently has only an infinitely small value as long as integration extends only from $0$ to an infinitely small value of $\sigma$. For any finite value of $\sigma$, however, the coefficient
\SourceLeaf{224}%
of $\mathrm d\sigma$ has the same value up to an infinitely small difference, whether we take $x=0$ or take it to be infinitely small. This therefore applies to the whole integral as well when it extends from $\sigma=0$ to $\sigma=\BareRoot[m]{\rho'}$.

Our conclusions cease to hold in just one case, namely when $\frac a\rho$ is no longer of the same order as any power of $\rho$, as would happen, for example, if $\frac a\rho$ were of the same order as $\frac{1}{\log\frac1\rho}$. In this case $Q$ would grow without bound as the point $O$ came infinitely close to the surface. The same would hold for $X$ if this behavior occurred for every value of $\theta$, rather than for only one or a few. There is no need to develop this further here, since we can exclude this singular case from our investigation altogether without loss.

\Article{17}
We will now use the same assumptions and notation as in Article 15 to examine the quantity $Y$, of which $\frac{hb\mathrm db.\mathrm dc}{r^3}$ is an arbitrary element. Since $r=\BareRoot{(bb+cc+(a-x)^2)}$, we have
\[
\frac{\mathrm d\frac hr}{\mathrm db}=-\frac{hb}{r^3}+\frac1r.\frac{\mathrm dh}{\mathrm db}-\frac{h(a-x)}{r^3}.\frac{\mathrm da}{\mathrm db}
\]
when $c$ is held constant. The first integration on that basis therefore gives
\[
\int\frac{hb\mathrm db}{r^3}=\frac{h^*}{r^*}-\frac{h^{**}}{r^{**}}+\int\frac1r.\frac{\mathrm dh}{\mathrm db}.\mathrm db-\int\frac{h(a-x)}{r^3}.\frac{\mathrm da}{\mathrm db}.\mathrm db
\]
where the integrations run from the smallest to the largest value of $b$ for each given value of $c$, and $h^*,r^*,h^{**},r^{**}$ denote the values of $h$ and $r$ at those bounds, If we write for brevity
\[
\frac{h^*}{r^*}-\frac{h^{**}}{r^{**}}=T,\quad\frac\rho r.\frac{\mathrm dh}{\mathrm db}-\frac{h(a-x)\rho}{r^3}.\frac{\mathrm da}{\mathrm db}=U
\]
then
\[
Y=\int T\mathrm dc+\iint\frac U\rho.\mathrm db.\mathrm dc
\]
\SourceLeaf{225}%
where the integration with respect to $c$ must extend from the smallest value of this coordinate on the surface to the largest. In the double integral, $\mathrm db.\mathrm dc$ represents the projection of an arbitrary element of the surface onto the plane of $b,c$, so we may also write $\rho\mathrm d\rho.\mathrm d\theta$ for it. We then obtain
\[
Y=\int T\mathrm dc+\iint U\mathrm d\rho.\mathrm d\theta
\]
where the double integral must be taken from $\rho=0$ to $\rho=\rho'$ and from $\theta=0$ to $\theta=2\pi$. Reasoning as in Article 15, we readily see that the values of this expression agree up to infinitely small differences whether we take $x=0$ or take it to be infinitely small. In other words, the value of $Y$ has the same limit as positive and negative values of $x$ become infinitely small, and this limit is exactly the value of the above formula when we set $x=0$. By analogy we will denote this value by $Y^0$. We must point out, however, that we cannot call this \textit{the} value of $\int\frac{kb\mathrm ds}{r^3}$ for $x=0$ (since this expression for $x=0$ cannot be integrated in the true sense). It is only \textit{a} value of the integral, the one obtained by integrating in the order used above.

This result needs a qualification, similar to the one in Article 16, in the singular case where the radii of curvature at $P$ are infinitely small, and likewise when $\frac{\mathrm dh}{\mathrm db}$ becomes infinitely large there. For our purpose, however, there is no need to consider these exceptional cases separately. They can occur only at individual points or along individual lines (thus only at the boundaries of parts of the surface, never within those parts).

Finally, it is evident that the quantity $Z$, or the integral $\int\frac{kc\mathrm ds}{r^3}$, behaves exactly as $Y$ does. When the point $O$ comes infinitely close to the point $P$ along the first coordinate axis, this integral has the same limit $Z^0$ whether the approach is from the positive or the negative side. This limit is also the value of $\iint\frac{hc\mathrm dc.\mathrm db}{r^3}$ for $x=0$ when we integrate first with respect to $c$.

\SourceLeaf{226}%


\Article{18}
Consider that the quantities $\frac{\mathrm dV}{\mathrm dx}$, $\frac{\mathrm dV}{\mathrm dy}$, $\frac{\mathrm dV}{\mathrm dz}$ agree without qualification with $X,Y,Z$ at every point of space off the surface itself, and that $V$ varies continuously everywhere. We can then readily infer from the results of the preceding article that at an infinitely small distance from $P$, or for infinitely small values of $x,y,z$, the value of $V$ is given, up to infinitely small quantities of higher order, by
\[
V^0+x(X^0-2\pi k^0)+yY^0+zZ^0
\]
when $x$ is positive, or by
\[
V^0+x(X^0+2\pi k^0)+yY^0+zZ^0
\]
when $x$ is negative. Here $V^0$ denotes the value of $V$ at the point $P$ itself, or for $x=0$, $y=0$, $z=0$. Let us therefore consider the values of $V$ along a straight line through $P$ that makes angles $A,B,C$ with the three coordinate axes. We denote an arbitrary segment of this line by $t$ and use $t^0$ for the value of $t$ at the point $P$.
If $t-t^0$ is infinitely small, then to within an infinitesimal of higher order we have
\[
V=V^0+(t-t^0)(X^0\cos A+Y^0\cos B+Z^0\cos C\mp2\pi k^0\cos A)
\]
with the upper sign for positive and the lower sign for negative values of $(t-t^0)\cos A$. In other words, $\frac{\mathrm dV}{\mathrm dt}$ at the point $P$ has two different values for an acute angle $A$, namely
\[
\begin{gathered}
X^0\cos A+Y^0\cos B+Z^0\cos C-2\pi k^0\cos A\text{ and}\\
X^0\cos A+Y^0\cos B+Z^0\cos C+2\pi k^0\cos A
\end{gathered}
\]
according to whether $\mathrm dt$ is taken as positive or negative. When $A$ is a right angle, so that the straight line only touches the surface, the two expressions coincide and we obtain
\[
\frac{\mathrm dV}{\mathrm dt}=Y^0\cos B+Z^0\cos C
\]
\begin{center}*\qquad\qquad*\\{}*\end{center}

\SourceLeaf{227}%
The propositions presented so far contain nothing essentially new, but could not be omitted for the sake of continuity. as necessary preliminaries to the investigations that follow, in. which a series of new theorems will be developed.

\Article{19}
Let $V$ be the potential of a system of masses $M',M'',M'''\ldots$ located at the points $P',P'',P'''\ldots$. Let $v$ be the potential of a second system of masses $m',m'',m'''\ldots$ placed at the points $p',p'',p'''\ldots$. Further, let $V',V'',V'''\ldots$ be the values of $V$ at the latter points, and $v',v'',v'''\ldots$ the values of $v$ at the points $P',P'',P'''\ldots$ We then have the equation
\[
M'v'+M''v''+M'''v'''+\text{etc.}=m'V'+m''V''+m'''V'''+\text{etc.}
\]
which can also be written as $\sum Mv=\sum mV$, where $M$ stands for an arbitrary mass of the first system and $m$ for an arbitrary mass of the second. In fact, each of $\sum Mv$ and $\sum mV$ is simply the sum of all the combinations $\frac{Mm}{\rho}$, where $\rho$ denotes the distance between the points containing the respective masses $M,m$.

If the masses of one or both systems are distributed continuously along lines, over surfaces or throughout regions in space, instead of being located at discrete points, the equation above still holds when the corresponding integral is substituted for the sum.

Suppose, for example, that the second system of masses is distributed on a surface so that the surface element $\mathrm ds$ carries the mass $k\mathrm ds$. Then $\sum Mv=\int kV\mathrm ds$. If the first system is likewise distributed on a surface, so that the surface element $\mathrm dS$ contains the mass $K\mathrm dS$, we obtain $\int Kv\mathrm dS=\int kV\mathrm ds$. In the latter case it is important to note that this equation still holds if the two surfaces coincide. For brevity, we will only sketch here the main steps of a rigorous justification of this extension of the proposition. It is not hard to show that the two integrals, when they refer to one and the same surface, are the limits of the integrals for two separate surfaces
\SourceLeaf{228}%
as the distance between them becomes infinitely small. For this purpose we need only take the two surfaces to be equal and parallel. This proof is immediately evident only when the given surface is such that the normals at all its points make acute angles with one straight line. A surface that fails to meet this condition (as every closed surface does) must first be divided into two or more parts that individually satisfy it, which makes it easy to reduce this case to the preceding one.

\Article{20}
Let us apply the theorem of the preceding article when the second system of masses is distributed with uniform density $k=1$ on a spherical surface of radius $=R$. The resulting potential $v$ is constant inside the sphere and $=4\pi R$. At every point outside the sphere whose distance from its center is $=r$, we have $v=\frac{4\pi RR}{r}$. This is equal to the potential produced at the center by a mass $4\pi RR$ placed at that point. On the surface of the sphere the two values of $v$ coincide. Thus if the first system of masses lies wholly inside the sphere, $\sum Mv$ equals the product of the total mass of this system and $4\pi R$. If that system of masses lies wholly outside the sphere, $\sum Mv$ equals the product of the potential of this mass at the center of the sphere and $4\pi RR$. Finally, if the first system of masses is distributed continuously on the sphere's surface, either expression can be used for $\int Kv\mathrm dS$. This gives the following theorem.

\textsc{Theorem.} Let $V$ denote the potential of a mass distributed in any manner, evaluated at an element of a spherical surface of radius $R$, with the element denoted by $\mathrm ds$. Integrating over the whole spherical surface gives
\[
\int V\mathrm ds=4\pi(RM^0+RRV^0)
\]
where $M^0$ denotes the entire mass inside the sphere and $V^0$ the potential of the exterior mass at the center of the sphere, with any masses continuously distributed on the surface of the sphere assigned at will to the exterior or interior masses.

\SourceLeaf{229}%
\Article{21}
\textsc{Theorem.} The potential $V$ of masses that all lie outside a connected region cannot have a constant value in one part of that region and at the same time a different value in another part of the same region.

\textit{Proof.} Suppose that the potential at every point of region $A$ is constant $=a$, and that at every point of another mass-free region adjoining $A$, denoted by $B$, it is (algebraically) greater than $a$. Construct a sphere with one part in $B$ and the rest, including its center, in $A$. Such a construction will always be possible. If $R$ is the radius of this sphere and $\mathrm ds$ an arbitrary element of its surface, then by the theorem of the preceding article $\int V\mathrm ds=4\pi RRa$ and $\int(V-a)\mathrm ds=0$. This is impossible because on the part of the surface that lies in $A$ we have $V-a=0$, and on the rest the difference is, by hypothesis, not $=0$ but positive.

In much the same way, we see that in a region adjoining $A$, it is impossible for $V$ to be less than $a$ at every point.

Clearly, however, at least one of these two cases would have to occur if our theorem were false.

The theorem includes the following two propositions.

I. If the region containing the masses surrounds a mass-free region like a shell, and the potential has a constant value in part of the enclosed region, then that value holds at every point of the entire enclosed region.

II. If the potential of masses contained in a finite region has a constant value in any part of the exterior region, then that value holds throughout the infinite exterior region.

It is also easy to see that in this second case the constant value of the potential can only be $0$. For let $M$ denote the total of all the masses if they all have the same sign, or, if they have different signs, the total of the positive or of the negative masses alone, according as the former or the latter predominate. Then the potential at a point whose distance from the nearest mass element is $=r$ is in any event, in abso%
\SourceLeaf{230}%
lute value, less than $\frac Mr$. This fraction can clearly become smaller than any assignable quantity in the exterior region.

\Article{22}
\textsc{Theorem.} Let $\mathrm ds$ be an element of a surface bounding a connected finite region, and let $P$ be the force that masses distributed in any manner exert at $\mathrm ds$ in the direction normal to the surface. An inward or outward force is taken as positive according as attracting or repelling masses are regarded as positive. Then the integral $\int P\mathrm ds$ over the whole surface is $=4\pi M+2\pi M'$, where $M$ denotes the total of the masses inside the region and $M'$ the total of those distributed continuously on the surface.

\textit{Proof.} Denote by $U\mathrm d\mu$ the part of $P$ due to the mass element $\mathrm d\mu$, by $r$ the distance of the element $\mathrm d\mu$ from $\mathrm ds$, and by $u$ the angle that the inward normal at $\mathrm ds$ makes with $r$. Then $U=\frac{\cos u}{rr}$. For any given $\mathrm d\mu$, a theorem proved in Article 6 of the \textit{Theoria Attractionis corporum sphaeroidicorum ellipticorum} gives $\int\frac{\cos u}{rr}.\mathrm ds=0$, $2\pi$ or $4\pi$, according as $\mathrm d\mu$ lies outside the region bounded by the surface, on the surface itself, or inside that region. Since $\int P\mathrm ds$ equals the sum of all the contributions $\mathrm d\mu.\int U\mathrm ds$, our theorem follows directly.

It remains to point out that the auxiliary theorem used here, as stated in the work just cited, needs modification in one particular case. The quantity $r$ denotes the distance of a \textit{given point} from the element $\mathrm ds$. When this point lies on the surface itself, the formula $\int\frac{\cos u}{rr}.\mathrm ds=2\pi$ is valid only if the continuity of the surface's curvature is undisturbed at that point. That continuity is broken when the point lies on an edge or at a corner. In that case $2\pi$ must be replaced by the area of the figure that all the straight lines issuing from the point and tangent to the surface cut out of a spherical surface of radius $1$ centered at the point. However, these exceptional cases concern only lines or points, and thus concern no \textit{parts} of the surface, only the bound%
\SourceLeaf{231}%
aries between parts. They therefore have no effect on the application of the auxiliary theorem here, as is clear.

\Article{23}
We draw a normal through each point of the surface and denote by $p$ the distance of an arbitrary point on the normal from its origin on the surface itself, taking the distance as positive on the inner side of the surface. The potential $V$ of the masses can be regarded as a function of $p$ and two other variables that distinguish the individual points of the surface in some way. The same holds for the partial derivative $\frac{\mathrm dV}{\mathrm dp}$, whose value we will consider here only for points on the surface itself, or for $p=0$. This derivative is identical to $P$ when masses are present only in the interior, only in the exterior, or in both, and no mass is distributed on the surface itself. In that case we have
\[
\int\frac{\mathrm dV}{\mathrm dp}.\mathrm ds=4\pi M
\]
If, on the other hand, the entire mass is distributed on the surface alone, so that the element $\mathrm ds$ contains the mass $k\mathrm ds$, then $\frac{\mathrm dV}{\mathrm dp}$ and $P$ no longer have the same meaning. The latter quantity evidently plays the same role here with respect to $p$ as $X^0$ does with respect to $x$ in Article 15. The derivative $\frac{\mathrm dV}{\mathrm dp}$, however, has two different values, namely $P-2\pi k$ and $P+2\pi k$, according as $\mathrm dp$ is taken as positive or negative. Now $\int k\mathrm ds$ is clearly equal to the entire mass $M'$ distributed on the surface, and by the theorem of the preceding article $\int P\mathrm ds=2\pi M'$. We therefore have
\[
\int\frac{\mathrm dV}{\mathrm dp}.\mathrm ds=0\quad\text{or}\quad\int\frac{\mathrm dV}{\mathrm dp}.\mathrm ds=4\pi M'
\]
according as $\frac{\mathrm dV}{\mathrm dp}$ is understood everywhere to mean the value on the inner or the outer side of the surface. In the first case the integral $\int\frac{\mathrm dV}{\mathrm dp}.\mathrm ds$ therefore behaves exactly as if the mass $M'$ belonged to the exterior, and in the second as if it belonged to the interior.

Thus, however the masses are distributed, the equation
\[
\int\frac{\mathrm dV}{\mathrm dp}.\mathrm ds=4\pi M
\]
\SourceLeaf{232}holds in general in the sense that $M$ denotes the mass contained in the interior. It is understood that any masses continuously distributed on the surface itself must be included in or excluded from the interior masses according as the value of $\frac{\mathrm dV}{\mathrm dp}$ for the outer or inner side has been chosen.

If, then, the interior contains no masses at all, and if $\frac{\mathrm dV}{\mathrm dp}$ is always understood to mean the value on the inner side, we have
\[
\int\frac{\mathrm dV}{\mathrm dp}.\mathrm ds=0
\]
\Article{24}
We use the same assumptions as at the end of the preceding article and denote the region under consideration by $T$. At an element $\mathrm dT$ of the region, let the total force due to masses outside the region or continuously distributed on its surface be denoted by $q$. We then have the following important result.

\textsc{Theorem.} We have
\[
\int V\frac{\mathrm dV}{\mathrm dp}.\mathrm ds=-\int qq\,\mathrm dT
\]
when the first integral is taken over the whole surface and the second throughout the whole region $T$.

\textit{Proof.} Introducing rectangular coordinates $x,y,z$, we first consider a straight line parallel to the $x$ axis and cutting the region $T$, so that $y,z$ have constant values. From the identity
\[
\frac{\mathrm d}{\mathrm dx}(V\frac{\mathrm dV}{\mathrm dx})=(\frac{\mathrm dV}{\mathrm dx})^2+V\frac{\mathrm d\mathrm dV}{\mathrm dx^2}
\]
it follows that the integral
\[
\int((\frac{\mathrm dV}{\mathrm dx})^2+V\frac{\mathrm d\mathrm dV}{\mathrm dx^2})\,\mathrm dx
\]
taken along the part of that straight line lying within $T$ is the difference between the two values of $V\frac{\mathrm dV}{\mathrm dx}$ at the initial and final points if the line cuts the boundary surface only twice. In general it is $=\sum\varepsilon V\frac{\mathrm dV}{\mathrm dx}$, where for $V\frac{\mathrm dV}{\mathrm dx}$ we take its individual values at the dif\SourceLeaf{233}ferent points of intersection, and where $\varepsilon$ is $=-1$ at the odd-numbered intersections (the first, third, and so on) and $=+1$ at the even-numbered ones. Now consider along this straight line the prismatic region having the rectangle $\mathrm dy.\mathrm dz$ as a cross section, and thus having $\mathrm dx.\mathrm dy.\mathrm dz$ as an element. The integral
\[
\int((\frac{\mathrm dV}{\mathrm dx})^2+V\frac{\mathrm d\mathrm dV}{\mathrm dx^2})\,\mathrm dT
\]
taken throughout the part of $T$ that lies in this prismatic region is then $=\sum\varepsilon V\frac{\mathrm dV}{\mathrm dx}.\mathrm dy.\mathrm dz$. This prism cuts out two pieces of the boundary surface, or in general an even number of pieces. Let each be denoted by $\mathrm ds$ and let $\xi$ be the angle between the $x$ axis and the inward normal to $\mathrm ds$. Then $\mathrm dy.\mathrm dz=\pm\cos\xi.\mathrm ds$, taking the upper sign at the odd-numbered intersections and the lower sign at the even-numbered ones. The above integral is therefore
\[
=-\sum V\frac{\mathrm dV}{\mathrm dx}.\cos\xi.\mathrm ds
\]
where the sum runs over all the surface elements concerned. If we now divide the whole region $T$ into such prismatic elements, the corresponding pieces of the surface also cover it completely. We thus have
\[
\int((\frac{\mathrm dV}{\mathrm dx})^2+V\frac{\mathrm d\mathrm dV}{\mathrm dx^2})\,\mathrm dT=-\int V\frac{\mathrm dV}{\mathrm dx}.\cos\xi.\mathrm ds
\]
with the first integration extending throughout the region $T$ and the second over the whole surface. Clearly, $\cos\xi$ equals the partial derivative $\frac{\mathrm dx}{\mathrm dp}$, with $p$ having the meaning fixed in Article 23, and with $x$ viewed as a function of $p$ and two other variables distinguishing the individual points of the surface. Consequently,
\[
\int((\frac{\mathrm dV}{\mathrm dx})^2+V\frac{\mathrm d\mathrm dV}{\mathrm dx^2}).\mathrm dT=-\int V\frac{\mathrm dV}{\mathrm dx}.\frac{\mathrm dx}{\mathrm dp}.\mathrm ds
\]
It is also clear that when the surface itself contains masses and $\frac{\mathrm dV}{\mathrm dx}$ therefore has two different values, the value meant here is always the one for the interior.

By entirely similar reasoning we find
\SourceLeaf{234}\[
\begin{gathered}
\int((\frac{\mathrm dV}{\mathrm dy})^2+V\frac{\mathrm d\mathrm dV}{\mathrm dy^2})\,\mathrm dT=-\int V\frac{\mathrm dV}{\mathrm dy}.\frac{\mathrm dy}{\mathrm dp}.\mathrm ds\\
\int((\frac{\mathrm dV}{\mathrm dz})^2+V\frac{\mathrm d\mathrm dV}{\mathrm dz^2})\,\mathrm dT=-\int V\frac{\mathrm dV}{\mathrm dz}.\frac{\mathrm dz}{\mathrm dp}.\mathrm ds
\end{gathered}
\]

If we now add these three equations and note that in the region $T$ we have
\[
\begin{gathered}
\frac{\mathrm d\mathrm dV}{\mathrm dx^2}+\frac{\mathrm d\mathrm dV}{\mathrm dy^2}+\frac{\mathrm d\mathrm dV}{\mathrm dz^2}=0\\
(\frac{\mathrm dV}{\mathrm dx})^2+(\frac{\mathrm dV}{\mathrm dy})^2+(\frac{\mathrm dV}{\mathrm dz})^2=qq
\end{gathered}
\]
and on the boundary surface
\[
\frac{\mathrm dV}{\mathrm dx}.\frac{\mathrm dx}{\mathrm dp}+\frac{\mathrm dV}{\mathrm dy}.\frac{\mathrm dy}{\mathrm dp}+\frac{\mathrm dV}{\mathrm dz}.\frac{\mathrm dz}{\mathrm dp}=\frac{\mathrm dV}{\mathrm dp}
\]
we obtain $\int qq\,\mathrm dT=-\int V.\frac{\mathrm dV}{\mathrm dp}.\mathrm ds$, which is our theorem itself. Using the final proposition of the preceding article as well, we can express it more generally as
\[
\int qq\,\mathrm dT=\int(A-V)\frac{\mathrm dV}{\mathrm dp}.\mathrm ds
\]
where $A$ denotes an arbitrary constant.

\Article{25}
\textsc{Theorem.} If, under the same assumptions as in the preceding article, the potential $V$ has the same value at all points on the boundary surface of the region $T$, this value also holds at all. points of the region itself, and the forces cancel completely throughout the region.

\textit{Proof.} If in the extended theorem of the preceding article we take $A$ to be the constant boundary value of the potential, then clearly $\int qq\,\mathrm dT=0$. Hence necessarily $q=0$ at every point of the region $T$, and therefore also $\frac{\mathrm dV}{\mathrm dx}=0$, $\frac{\mathrm dV}{\mathrm dy}=0$, $\frac{\mathrm dV}{\mathrm dz}=0$. It follows that $V$ is constant throughout the region $T$.

\Article{26}
\textsc{Theorem.} Suppose the masses lie solely within the finite region $T$ or are also distributed continuously, wholly or in part, over i%
\SourceLeaf{235}%
ts surface $S$. If their potential has at every point of $S$ the same value $=A$, then at every point $O$ of the infinite exterior region $T'$ the potential will be

\noindent first, if $A=0$, likewise $=0$,

\noindent second, if $A$ is not $=0$, less than $A$ and of the same sign as $A$.

\textit{Proof.} I. First it must be shown that the potential at $O$ cannot take a value outside the bounds $0$ and $A$. Suppose the potential at $O$ has such a value $B$, and let $C$ denote any quantity that lies both between $B$ and $0$ and between $B$ and $A$. If straight lines are drawn from $O$ in every direction, each will contain a point $O'$ where the potential is $=C$, with the entire line $OO'$ lying in the region $T'$. This follows directly from the fact that the potential varies continuously. If the straight line is extended far enough, the potential either passes from $B$ to $A$ or diminishes indefinitely, depending on whether the line meets the surface $S$ (see the remark at the end of Article 21). The set of all points $O'$ then forms a closed surface. Since the potential on it is constant $=C$, the theorem of the preceding article requires that it have the same value at every point in the region enclosed by this surface. Yet at $O$ it has a value different from $C$, namely $B$. The assumption therefore necessarily leads to a contradiction.

This completes the proof of our theorem for $A=0$. In the second case, where $A$ is not $=0$, it shows that at no point in $T'$ can the potential be greater than $A$ or have the opposite sign.

II. To complete our proof in the second case, we draw a spherical surface centered at $O$ with radius $R$ less than the shortest distance from the point $O$ to $S$. We divide it into elements $\mathrm ds$ and denote the potential at each element by $V$. The potential at $O$ is again denoted by $B$. By the theorem of Article 20, the integral taken over the entire spherical surface is then
\[
\int V\mathrm ds=4\pi RRB,\quad\text{and hence}\quad\int(V-B)\mathrm ds=0
\]
\SourceLeaf{236}This equality can hold only if either $V$ is constant $=B$ at every point on the spherical surface, or $V$ differs from $B$ in opposite senses on different parts of the surface. Under the first assumption, by Art.\ 25 the potential would have to be constant throughout the interior of the sphere, and hence by Art 21 throughout the infinite region $T'$, with value $=0$. This conflicts with the assumption that at the boundary of this region, on the surface $S$, the potential differs from $0$ and with the impossibility of its changing abruptly from there. The second assumption, however, would conflict with what was proved under I: if $B$ were either $=0$ or $=A$. Therefore $B$ must lie \textit{between} $0$ and $A$.

\Article{27}
\textsc{Theorem.} In the theorem of the preceding article, the first case, corresponding to the value $0$ of the constant potential $A$, can occur only if the sum of all the masses is itself $=0$. The second case can occur only if this sum is not $=0$.

\textit{Proof.} Let $\mathrm ds$ be an element of the surface of any sphere enclosing the region $T$, $R$ its radius, $M$ the sum of all the masses, and $V$ their potential at $\mathrm ds$. By the theorem of Article 20 the integral is $\int V\mathrm ds=4\pi RM$. By the preceding theorem, in the first case, where $A=0$, the potential $V$ is $=0$ at every point of the spherical surface. In the second case it is less than $A$ and has the same sign. Thus in the first case $4\pi RM=0$ and hence $M=0$, whereas in the second case $4\pi RM$, and therefore $M$ as well, must have the same sign as $A$. It also follows that in the second case $4\pi RM$ is less than $\int A\mathrm ds$, or $4\pi RRA$, so $M$ is less than $RA$, or $A$ is greater than $\frac{M}{R}$.

The second part of this theorem, taken together with the theorem of the preceding article, can clearly also be stated as follows.

Suppose the algebraic sum of masses contained in a region bounded by a closed surface, or also distributed in part continuously on the surface itself, is $=0$, and their potential has a constant value at every point on the surface. That value must itself be $=0$
\SourceLeaf{237}and hold throughout the infinite exterior region as well. Consequently the effects of the forces due to those masses cancel completely throughout the exterior region.

\Article{28}
It is easy to see that all the conclusions of the two preceding articles remain valid if $S$ is a surface that is not closed and the masses lie only on that surface. Here the region $T$ drops out altogether. All points outside the surface itself belong to the infinite exterior region. If the potential has a constant value different from $0$, namely $A$, everywhere on the surface, then it has a smaller value of the same sign everywhere outside it.

What was said about the first case, $A=0$, remains true here too, but is vacuous. In this case the potential $V$ is $=0$ at every point in space. Hence $\frac{\mathrm dV}{\mathrm dt}=0$ everywhere, where $t$ denotes any straight line. From this and Article 18 it follows readily that the density on the surface must be $=0$ everywhere, so the surface cannot contain any mass at all.

This last remark applies generally whenever the masses are to lie only on the surface itself, even if it is closed. For by the theorem of Article 25 the value of the potential in this case will clearly be $=0$ throughout the interior region as well.

\Article{29}
Before we proceed to the following investigations, in which masses distributed continuously over a surface play a central role, we must point out an important distinction between distributions. Either we allow only masses of one sign (which for brevity we will always take to be positive), or we allow masses of opposite signs as well. Suppose a mass $M$ is distributed over a surface in such a way that each surface element $\mathrm ds$ carries the mass $m\mathrm ds$. As before, we call $m$ the density, and $\int m\mathrm ds$ taken over the whole surface is $=M$. We call this a \textit{single-sign} distribution if $m$ is positive everywhere, or at least negative nowhere. If, on the other hand, $m$ is positive in some places and negative in others, we call the distribu%
\SourceLeaf{238}%
tion \textit{mixed-sign}. In that case $M$ is only the algebraic sum of the portions of mass, or the absolute difference between the positive and negative masses. A very special case of a mixed-sign distribution occurs when $M=0$. It may then admittedly seem awkward to go on speaking of the mass $0$ as distributed over the surface.

\Article{30}
However a mass $M$ is spread over a surface in a \textit{single-sign} distribution, it is evident that the resulting potential $V$ is everywhere positive and will be greater at every point of the surface than $\frac{M}{r}$, where $r$ denotes the greatest distance between two points of the surface. The potential could attain that very value only at one endpoint of the line $r$ if all the mass were concentrated at the other endpoint. That case does not arise here, since we are concerned only with continuous distribution, in which each surface element $\mathrm ds$ has only an infinitesimal mass $m\mathrm ds$. The integral $\int Vm\mathrm ds$ taken over the whole surface is therefore in any case greater than $\int\frac{M}{r}m\mathrm ds$, or $\frac{MM}{r}$, and so there must be a single-sign distribution for which this integral attains a minimum. We may state in advance that one aim of the following investigations is to prove that, for such a distribution, in which $\int Vm\mathrm ds$ attains its minimum, the potential $V$ will have the same value at every point of the surface, that no part of the surface can remain unoccupied, and that there is only one such distribution. For brevity, however, we will conduct the investigation in a more general form from the outset.

\Article{31}
Let $U$ be a quantity with a definite finite value at every point on the surface, varying continuously. The integral
\[
\Omega=\int(V-2U)m\mathrm ds
\]
taken over the whole surface can then have very different values for different single-sign distributions of the mass $M$. Clearly, however,
\SourceLeaf{239}there must be one such distribution for which this integral attains a minimum. A proof will now be given of the following

\textsc{Theorem.} For such a distribution,

1. the difference $V-U=W$ has a constant value everywhere on the surface where it is actually occupied by portions of $M$,

2. if any parts of the surface remain unoccupied, $W$ must be greater there, or at least cannot be smaller, than that constant value.

I. First it must be shown that if one distribution is replaced by another differing from it by an infinitely small amount, by substituting $m+\mu$ for $m$, the resulting variation of $\Omega$ is expressed by $2\int W\mu\mathrm ds$.

Indeed, if we denote the variations of $\Omega$ and $V$ by $\delta\Omega$ and $\delta V$, then
\[
\delta\Omega=\int\delta V.m\mathrm ds+\int(V-2U)\mu\mathrm ds
\]
But at the same time $\int\delta V.m\mathrm ds=\int V\mu\mathrm ds$, as follows readily from the theorem of Article 19. For $\delta V$ is simply the potential of the mass distribution whose density in each surface element is $\mu$. Thus what is denoted here by $V,m,\delta V,\mu$ can be substituted there for $V,K,v,k$, with $\mathrm ds$ standing for both $\mathrm dS$ and $\mathrm ds$. Consequently,
\[
\delta\Omega=\int(2V-2U)\mu\mathrm ds=2\int W\mu\mathrm ds
\]

II. The variations $\mu$ must clearly satisfy the general condition $\int\mu\mathrm ds=0$. For the present investigation they must also satisfy a second condition. The value of $\mu$ cannot be negative in the unoccupied parts of the surface, if there are any, since otherwise the distribution would cease to be single-sign.

III. Suppose now that, for a given distribution of $M$, the quantity $W$ has unequal values in different parts of the surface. Let $A$ be a quantity between the unequal values of $W$. Let $P$ be the part of the surface where the values of $W$ are greater, and $Q$ the part where they are smaller, than $A$. Further, let $p,q$ be parts of the surface of equal area, the former belonging to $P$ and the latter to $Q$. Under these assumptions, we give the variation
\SourceLeaf{240}of $m$ everywhere in $p$ the constant negative value $\mu=-\nu$, everywhere in $q$ the positive value $\mu=\nu$, and the value $0$ in all the remaining parts of the surface. This clearly satisfies the first condition in II. The second condition further requires that $p$ contain no unoccupied parts. This can always be achieved, provided the whole of $P$ is not unoccupied.

The result, however, is that $\delta\Omega$ takes a negative value, as is easy to see by writing this variation in the form $2\int(W-A)\mu\mathrm ds$.

It follows that, if a given distribution has unequal values of $W$ in the occupied part of the surface, or if the values in the occupied part are equal but smaller values occur in the unoccupied part, then changing the distribution can reduce $\Omega$. At the minimum, therefore, the conditions stated in the theorem above must hold.

\Article{32}
For our special case (Art.\ 30), where $U=0$, so that $W$ is simply the potential of the mass distributed over the surface and $\Omega$ denotes the integral $\int Vm\mathrm ds$, let us now combine the theorem of the preceding article with the one cited in Article 28. It follows immediately that when $\int Vm\mathrm ds$ attains its minimum, the surface can have no unoccupied parts. Otherwise, even if the whole surface is closed, the occupied part would have to be regarded as an unclosed surface, and the unoccupied part as belonging to the exterior region relative to it. By Art.\ 28 the potential there would therefore have to be smaller than on the occupied surface, whereas the theorem of the preceding article rules out a smaller value.

It is thus proved that a given mass has a single-sign distribution over the whole surface in which no part remains empty and which produces the same potential at every point on the surface. The only remaining step needed to complete the proof of the theorem stated in Article 30 is to show that only one distribution can do this. That will follow below as part of a more general theorem.

The requirement that no part of the surface remain unoccupied if $\int Vm\mathrm ds$ is to attain its minimum can clearly also be expressed as follows.
\SourceLeaf{241}For every distribution that leaves a finite part of the surface empty, the integral $\int Vm\mathrm ds$ exceeds its minimum by a finite difference.

\Article{33}
The vital nerve of the argument developed in Article 31 lies in how immediately evident it is that a minimum value of $\Omega$ exists as long as we restrict ourselves to single-sign distributions of a given mass. If this were equally evident without that restriction, the conclusions there would directly yield the result \textit{that there is always, if not a single-sign distribution, at least a mixed-sign distribution of the given mass, for which $W=V-U$ has a constant value at every point on the surface}, since the second condition (Art.\ 31.\ II) would then drop out. But since this ceases to be evident as soon as we remove the restriction to single-sign distributions, we must seek a rigorous proof of this most important theorem of our entire investigation by a somewhat more contrived route. The following seems the simplest way to reach the goal.

We first consider three different mass distributions, for which we use the following particular symbols in place of the general symbols for density $m$ and potential $V$.
\[
\begin{array}{rll}
\mathrm{I.}&m=m^0,&V=V^0\\
\mathrm{II.}&m=m',&V=V'\\
\mathrm{III.}&m=\mu,&V=v
\end{array}
\]
Distribution I is the single-sign distribution of the positive mass $M$ for which $\int Vm\mathrm ds$ attains its minimum.

II. is the single-sign distribution of the same mass $M$ for which $\int(V-2\varepsilon U)m\mathrm ds$ attains its minimum, where $\varepsilon$ is any constant coefficient.

III. depends on I and II through $\mu=\frac{m'-m^0}{\varepsilon}$, and is alo a mixed-sign distribution in which the total mass is $=0$.

By what was proved in Article 31, $V^0$ is constant over the entire
\SourceLeaf{242}surface; $V'-\varepsilon U$ is constant on the surface as far as it is occupied in the second distribution, and hence $v-U$ is also constant on that same part of the surface, since $v=\frac{V'-V^0}{\varepsilon}$.

Whether the second distribution occupies the whole surface, or leaves a larger or smaller part unoccupied, depends on the coefficient $\varepsilon$. Since the second distribution becomes the first when $\varepsilon=0$, generally speaking, the part of the surface left unoccupied for a given value of $\varepsilon$ will shrink as $\varepsilon$ decreases and will become entirely filled before $\varepsilon$ reaches the value $0$. In exceptional cases, however, a part may always remain unoccupied as long as $\varepsilon$ differs from $0$ and does not take the opposite sign. For our purpose it is enough to take $\varepsilon$ infinitely small. It is then easy to show that no finite part of the surface can in any case remain unoccupied. Otherwise, by the concluding remark of Art.\ 32, the integral $\int V'm'\mathrm ds$ would have to exceed $\int V^0m^0\mathrm ds$ by a finite difference. If this difference is denoted by $e$, the difference between the two integrals is
\[
\int(V'-2\varepsilon U)m'\mathrm ds-\int(V^0-2\varepsilon U)m^0\mathrm ds=e-2\varepsilon\int U(m'-m^0)\mathrm ds
\]
which remains positive for an infinitely small $\varepsilon$. This contradicts the assumption that $\int(V-2\varepsilon U)m\mathrm ds$ attains its minimum in the second distribution.

It follows that, if in the third distribution we take for $\mu$ the limiting value of $\frac{m'-m^0}{\varepsilon}$ as $\varepsilon$ decreases indefinitely, then $v-U$ has a constant value over the whole surface.

If we now form a fourth distribution by setting $m=m^0+\mu$, so that the total mass remains $=M$, the resulting potential will be $=V^0+v$. It will therefore exceed $U$ over the entire surface by the constant difference $V^0+v-U$. This proves the theorem stated above.

\Article{34}
It remains to prove that only one distribution of a given mass $M$ is possible for which $V-U$ is constant over the entire surface. Suppose there were two different distribu%
\SourceLeaf{243}%
tions with this property. Denote $m$ and $V$ in the first by $m'$, $V'$, and in the second by $m''$, $V''$. Then a third mass distribution, in which $m=m'-m''$ is taken, would have potential $=V'-V''$, which is therefore constant, and total mass $=0$. By Art.\ 27 this constant potential must be $=0$. Consequently, by Art.\ 28, $m'-m''=0$ as well, so the two distributions are identical.

Finally, it must be mentioned that there is always a mass distribution for which the difference $V-U$ has a \textit{given} constant value. Let $\alpha$ be any constant coefficient, and retain the notation for the first and third distributions in the preceding article. The potential of the distribution in which $m=\alpha m^0+\mu$ is taken will be $=\alpha V^0+v$, and the constant difference $\alpha V^0+v-U$ can be given any value by choosing the coefficient $\alpha$ suitably. The total mass of this distribution is then no longer arbitrary, but is $=\alpha M$. The same reasoning as before shows that this condition on the distribution can also be fulfilled in only one way.

\Article{35}
In most cases, actually determining the distribution of mass on a given surface for every prescribed form of $U$ exceeds the powers of analysis in its present state. The simplest case within our reach is that of an entire spherical surface. We will, however, proceed directly to the more general case where the surface differs very little from a spherical surface and quantities of higher order than the deviation itself may be neglected.

Let $R$ be the radius of the sphere, $r$ the distance of any point in space from its center, $u$ the angle between $r$ and a fixed straight line, and $\lambda$ the angle between the plane through that straight line and $r$ and a fixed plane. Let the distance of an arbitrary point on the given closed surface from the center of the sphere be $=R(1+\gamma z)$, where $\gamma$ is a very small constant factor whose higher powers are neglected, while $z$, like $U$, is a function of $u$ and $\lambda$.

\SourceLeaf{244}The potential $V$ of the mass distributed over the spherical surface is expressed at every point of the exterior region by a series in descending powers of $r$, to which we give the form
\[
A^0\frac Rr+A'\left(\frac Rr\right)^2+A''\left(\frac Rr\right)^3+\text{etc.}
\]
and at every point of the interior region by the ascending series
\[
B^0+B'\frac rR+B''\left(\frac rR\right)^2+B'''\left(\frac rR\right)^3+\text{etc.}
\]
The coefficients $A^0,A',A''$ and so on are functions of $u$ and $\lambda$ that satisfy known partial differential equations, \textit{Allg.\ Th.\ d.\ Erdm.}\ Art.\ 18, as do $B^0,B',B''$ and so on. On the given surface, the potential is to equal a given function of $u$ and $\lambda$, namely $V=U$, so
\[
\left(\frac rR\right)^{\frac12}V=(1+\gamma z)^{\frac12}U
\]
Suppose, then, that $(1+\gamma z)^{\frac12}U$ has been expanded in a series
\[
P^0+P'+P''+P'''+\text{and so on}
\]
such that the individual terms $P^0$, $P'$, $P''$, $P'''$ and so on also satisfy those differential equations. Since both series for the potential above must remain valid up to the surface itself, it follows that
\[
\begin{aligned}
&P^0+P'+P''+P'''+\text{etc.}\\
&\qquad=A^0(1+\gamma z)^{-\frac12}+A'(1+\gamma z)^{-\frac32}+A''(1+\gamma z)^{-\frac52}+\text{etc.}\\
&\qquad=B^0(1+\gamma z)^{\frac12}+B'(1+\gamma z)^{\frac32}+B''(1+\gamma z)^{\frac52}+\text{etc.}
\end{aligned}
\]
From this we conclude that, if quantities of order $\gamma$ are neglected,
\[
P^0+P'+P''+\text{etc.}=A^0+A'+A''+\text{etc.}
\]
and hence (since a function of $u$, $\lambda$ can be expanded in only one way in a series whose terms satisfy the differential equations mentioned)
\[
P^0=A^0,\qquad P'=A',\qquad P''=A''\quad\text{etc.}
\]
Similarly, neglecting quantities of order $\gamma$,
\[
P^0=B^0,\qquad P'=B',\qquad P''=B''\quad\text{etc.}
\]

\SourceLeaf{245}If we therefore set (I)
\[
\begin{aligned}
A^0&=P^0+\gamma a^0,& B^0&=P^0-\gamma b^0\\
A'&=P'+\gamma a',& B'&=P'-\gamma b'\\
A''&=P''+\gamma a'',& B''&=P''-\gamma b''\\
A'''&=P'''+\gamma a''',& B'''&=P'''-\gamma b'''\\
&\text{etc.}
\end{aligned}
\]
where clearly $a^0,a',a'',a'''$ and so on, and likewise $b^0,b',b'',b'''$ and so on, also satisfy the differential equations mentioned, and substitute these values into the equations above while neglecting quantities of order $\gamma\gamma$, then, after division by $\gamma$, we obtain, up to errors of order $\gamma$,
\[
\begin{aligned}
a^0+a'+a''+a'''+\text{etc.}&=\tfrac12 z(P^0+3P'+5P''+7P'''+\text{etc.})\\
b^0+b'+b''+b'''+\text{etc.}&=\tfrac12 z(P^0+3P'+5P''+7P'''+\text{etc.})
\end{aligned}
\]
Thus, up to errors of order $\gamma$,
\[
b^0=a^0,\qquad b'=a',\qquad b''=a''\quad\text{etc.}
\]
and consequently, up to errors of order $\gamma\gamma$, (II)
\[
B^0=P^0-\gamma a^0,\qquad B'=P'-\gamma a',\qquad B''=P''-\gamma a''\quad\text{etc.}
\]

The derivative $\frac{\mathrm dV}{\mathrm dr}$ has two different values on the surface itself. The value corresponding to a negative $\mathrm dr$, or to the inner side, exceeds the one on the outer side by $4\pi m\cos\theta$, where $m$ is the density at the point of intersection and $\theta$ the angle between $r$ and the normal (Art.\ 13, where $t,A,k^0$ mean the same as $r,\theta,m$ here). These two values are found by differentiating the two expressions for $V$ valid in the interior and exterior regions with respect to $r$, and then setting $r=R(1+\gamma z)$. Thus the first is
\[
=\frac1R\{B'+2B''(1+\gamma z)+3B'''(1+\gamma z)^2+\text{etc.}\}
\]
and the second
\[
-\frac1R\{A^0(1+\gamma z)^{-2}+2A'(1+\gamma z)^{-3}+3A''(1+\gamma z)^{-4}+\text{etc.}\}
\]

\SourceLeaf{246}Multiplying the difference by $R(1+\gamma z)^{\frac32}$, we therefore have
\[
\begin{aligned}
&4\pi mR\cos\theta.(1+\gamma z)^{\frac32}\\
&\quad=\phantom{{}+{}}A^0(1+\gamma z)^{-\frac12}+2A'(1+\gamma z)^{-\frac32}+3A''(1+\gamma z)^{-\frac52}+\text{etc.}\\
&\qquad+B'(1+\gamma z)^{\frac32}+2B''(1+\gamma z)^{\frac52}+3B'''(1+\gamma z)^{\frac72}+\text{etc.}
\end{aligned}
\]
Substituting the values from I for $A^0,A'$ and so on, and those from II for $B^0,B'$ and so forth, and omitting quantities of order $\gamma\gamma$, we obtain
\[
\begin{aligned}
4\pi mR\cos\theta.(1+\gamma z)^{\frac32}&=P^0+3P'+5P''+7P'''+\text{etc.}\\
&\quad+\gamma(a^0+a'+a''+a'''+\text{etc.})\\
&\quad-\tfrac12\gamma z(P^0+3P'+5P''+\text{etc.})
\end{aligned}
\]
and consequently, since the last two series cancel one another up to quantities of order $\gamma\gamma$,
\[
m=\frac{(1+\gamma z)^{-\frac32}}{4\pi R\cos\theta}.(P^0+3P'+5P''+7P'''+\text{etc.})
\]
which solves the problem. Instead of $(1+\gamma z)^{-\frac32}$ we can also write $1-\frac32\gamma z$ and omit the divisor $\cos\theta$, insofar as, at least generally speaking, $\theta$ is of order $\gamma$ and therefore $\cos\theta$ differs from $1$ only by a quantity of order $\gamma\gamma$.

For a sphere, where $\gamma=0$, we have exactly
\[
m=\tfrac1{4\pi R}(P^0+3P'+5P''+7P'''+\text{etc.})
\]
where $P^0+P'+P''+P'''+\text{etc.}$ is the expansion of $U$ itself.

\Article{36}
The quantity $U$ has been left unspecified in the investigations so far. Applying them to the case where $U$ is the potential of a given system of masses now leads us to the following important

\textsc{Theorem.} Any given mass distribution $D$ confined either entirely to the interior region bounded by a closed surface $S$ or entirely to the exterior region can be replaced by a mass distribution $E$ confined to the surface itself, so that the
\SourceLeaf{247}effect of $E$ equals the effect of $D$ at every point of the exterior region in the first case, or at every point of the interior region in the second.

The only requirement is the following. If the potential of $D$ at every point of $S$ is denoted by $U$ and the potential of $E$ by $V$, then over the entire surface $V-U=0$ must hold in the first case, while in the second case the difference need only be constant. For $-U$ is the potential of a distribution $D'$ opposite to $D$ (each portion of mass being replaced by its opposite). Thus $V-U$ is the potential of the distributions $D'$ and $E$ taken together. Their effects therefore cancel throughout the exterior region in the first case, and throughout the interior region in the second (Artt.\ 27 and 25). In other words, the effects of $D$ and $E$ are equal in the respective regions. Moreover, in the first case the total mass in $E$ equals that in $D$, whereas in the second it remains arbitrary.

The theorem announced in the \textit{Intensitas vis magneticae}, Art.\ 2, and also cited at various points in the \textit{Allgemeine Theorie des Erdmagnetismus}, now appears as a special case of what has been proved here.

\Article{37}
Although, as already noted in Article 35, actually finding the complete distribution $E$ presents insurmountable difficulties in most cases, there is one case where it can be carried out very easily and deserves particular mention here. This is the case where $U$ is constant, so that $S$ is an equilibrium surface for the given system of masses $D$. It is easy to see that only the case need be considered where $D$ lies in the interior region and the total mass is not $=0$. Otherwise there would be no effect to replace by a mass distribution on $S$.

Let $O$ be a point of the surface $S$, and $r$ a straight line that cuts the surface there at right angles and is regarded as increasing from the inside toward the outside. Further, let $-C$ be the value of the derivative $\frac{\mathrm dU}{\mathrm dr}$ at $O$, and let $m$ be the density of the mass distribution $E$ at $O$. The derivative $\frac{\mathrm dV}{\mathrm dr}$ has two different values at $O$. The value on the outer side,
\SourceLeaf{248}since $V=U$ on the surface and throughout the exterior region, equals the derivative $\frac{\mathrm dU}{\mathrm dr}$ and is therefore $=-C$. The value on the inner side, however, is $=0$, since $V$ is constant on the surface and throughout the interior region. But the second value exceeds the first by $4\pi m$, so we have $4\pi m=C$ or $m=\frac C{4\pi}$. Clearly, $C$ is simply the force arising from the mass distribution $D$ and has the same sign as the total mass.

